\documentclass[12pt]{article}
\usepackage[utf8]{inputenc}
\usepackage[T1]{fontenc}
\usepackage{hyperref}
\title{Tigak Language Resources for K-12 Curriculum Development}
\author{Compiled by Claude Sonnet 5 (claude-sonnet-5) for Andr\'as Kornai}
\date{}
\begin{document}
\maketitle

\section{Introduction}
Tigak (also called Omo; ISO 639-3 tgc; Glottocode tiga1245) is an Austronesian language of the Kavieng District, New Ireland Province, Papua New Guinea. Wikipedia cites about 6,000 speakers (a 1991 figure). The search found one substantial description, a 1974 thesis, and no school materials.

\section{Primary Educational Resources}
No primers, textbooks or school programmes were found in the search results, although the reference description says Tigak is used for instruction in some schools. This could not be confirmed.

\section{Linguistic and Reference Materials}
\begin{itemize}
\item Clive Beaumont, \emph{The Tigak Language of New Ireland}, ANU PhD thesis, 1974, a tagmemic description of grammar and phonology: \url{https://openresearch-repository.anu.edu.au/bitstream/1885/130442/2/b1015637_Beaumont_1974_PhD.pdf}
\item Wikipedia overview of the Tigak language: \url{https://en.wikipedia.org/wiki/Tigak_language}
\item ANU record, ``New Ireland languages: a review'': \url{https://digitalcollections.anu.edu.au/items/0c80f500-3894-43b0-80e1-2201ff00e6bb/full}
\item ANU repository item (title ``Cultural advice''; content not verified): \url{https://openresearch-repository.anu.edu.au/items/3af0be78-62ba-4acf-9ba3-dffc4b0eeb30}
\item Digital Pasifik item (contents not verified): \url{https://digitalpasifik.org/items/356599}
\end{itemize}
Sources disagree on the phoneme inventory: one ANU record gives 12 consonants and 5 vowels, while Beaumont's 1969 paper gives 13 consonants and 8 vowels (three of them glides).

\section{Religious and Community Resources}
None found in the search results.

\section{Digital and Audio Resources}
No audio, video or software were found.

\section{Curriculum Development Context}
Little exists for curriculum use. The 1974 thesis and the 1969 phonology paper are the main basis for orthography and grammar work, and community partners would need to be identified locally. The claim of use in some schools needs on-site verification.
\end{document}